\documentclass[10pt, letterpaper, twocolumn]{article}

% Core packages
\usepackage[utf8]{inputenc}
\usepackage[T1]{fontenc}
\usepackage{mathpazo}
\usepackage{microtype}
\usepackage[margin=0.8in, columnsep=0.25in]{geometry}
\usepackage{setspace}
\singlespacing

% Formatting packages
\usepackage{titlesec}
\usepackage{fancyhdr}
\usepackage{booktabs}
\usepackage{enumitem}
\usepackage{hyperref}
\usepackage{xcolor}
\usepackage{tabularx}
\usepackage{threeparttable}

% Custom section styling
\renewcommand{\thesection}{\Roman{section}}
\renewcommand{\thesubsection}{\arabic{section}.\arabic{subsection}}

\titleformat{\section}
  {\normalfont\large\bfseries\scshape\centering}{\thesection.}{0.5em}{}
\titleformat{\subsection}
  {\normalfont\normalsize\bfseries}{\thesubsection}{0.5em}{}

% Header/Footer setup
\pagestyle{fancy}
\fancyhf{}
\fancyhead[L]{\small \scshape Autonomy and the Right to Check Out Permanently}
\fancyhead[R]{\small \thepage}
\renewcommand{\headrulewidth}{0.5pt}
\renewcommand{\footrulewidth}{0pt}

% Metadata
\title{\textbf{\LARGE AUTONOMY AND THE RIGHT TO\\ CHECK OUT PERMANENTLY}\\[0.5em]
\large \textit{Or: Why I Get to Champion Suicide Rights\\Now That I No Longer Want to Eat a Bullet}}
\author{\textbf{Justin Bogner}}
\date{2026}

\begin{document}

% --- Title and Author's Note (Full Width) ---
\twocolumn[
  \begin{@twocolumnfalse}
    \maketitle
    \vspace{-1em}
    \begin{center}
        \rule{0.6\textwidth}{0.5pt}
        \vspace{1em}

        \begin{minipage}{0.8\textwidth}
            \small
            \textbf{\scshape Author's Note}\\
            \itshape
            \setlength{\parindent}{1em}
            This essay intentionally blends personal reflection with a formal academic structure. The subject matter is intensely personal, but the argument it forwards is a serious political and ethical one, deserving of a rigorous presentation.

            The perspective here is not one of crisis, but of long-term reflection from a position of stability. It is an argument for bodily sovereignty, made after two decades of considering the experience from every angle. The language, including its profanities, is chosen for precision.
        \end{minipage}
        \vspace{1em}
        \rule{0.6\textwidth}{0.5pt}
        \vspace{2em}
    \end{center}
  \end{@twocolumnfalse}
]

% --- Abstract ---
\section*{Abstract}

I was doom-scrolling MAiD proposals---Canada's expanded eligibility debates---when something landed: I am not suicidal anymore. Not a Hallmark epiphany. A flat, quiet \textit{huh}. And the silence where that low-grade hum used to live was so unfamiliar it took a moment to name.

This essay proceeds from that silence. It is not a recovery narrative. It is a political argument about decisional capacity as the only legitimate gate for Medical Aid in Dying eligibility---argued from a position of stability, not suffering. The claim is simple: if my life is mine when I am happy and choosing, it was mine when I was finished and done. The state's discomfort with that conclusion does not constitute a counterargument. Capacity equals sovereignty. The rest is noise.

% --- Section I ---
\section{The Silence}

\subsection{What Stopped Humming}

For years---the better part of two decades---there was a sound. Not audible. More like the physical register of a refrigerator compressor: present so constantly that you stop hearing it, and only notice it when it stops. Mine was the low, continuous frequency of \textit{done-ness}. The background knowledge that exit was on the table. Always. Even during genuinely good stretches: cuddling the dog, watching the farm do what farms do, grinding out another book, running through the standard-issue Reasons to Stick Around. The hum was still there. I'd tuned to it the way you tune to traffic noise.

Then it wasn't.

I noticed because I was reading about Canada's Track 2 MAiD provisions---the non-terminal eligibility pathway, where the numbers climbed to 732 provisions in 2024, where 61.5\% of recipients report disability as a primary driver---and I thought, idly and without drama: \textit{I actually wouldn't do that}. Not ``I can't.'' Not ``I shouldn't.'' Not the guilt-trip reflexes that were installed in a psych ward when I was eighteen. Just: \textit{wouldn't}. Choice, native and quiet. The option still present. The desire simply gone.

I sat with it for a long time, suspicious of myself. The self-suspicion is trained in, too---years of being told that my judgment about my own interior states was not to be trusted, that the fact I thought I was fine was itself evidence I wasn't. Eventually I laughed. Ugly, genuine, a little undignified. Then something that might be called gratitude arrived: not for being saved, not for the intervention, but for the specific unexpected gift of a morning where I wasn't quietly wrestling with existence before the coffee finished brewing.

And then---right on the heels of that warm fuzzy---the political reflex, crisp as ever: this doesn't change anything. My life was mine when I was done. It is mine now that I'm not. The logic doesn't move with the mood.

\subsection{Why Stability Makes the Argument Stronger}

I know how the objection runs. \textit{Kid who tried to off himself at eighteen now advocates for suicide rights. Unfinished business. Philosopher cosplay. Crisis talking through a logic costume.}

The objection has it exactly backward.

Arguments for exit rights made from within crisis are epistemically suspect---not because they're wrong, but because they're impossible to cleanly separate from the state that generated them. The person in acute agony advocating for their own death can always be told their judgment is distorted, their preferences are symptoms, their desire for exit is a thing to be treated rather than heard. This is the rhetorical trap the protectionist system relies on: any statement of preference becomes evidence of incapacity.

Arguing from stability closes that trap. I am not rationalizing a plan. I have no plan. I'm in what an assessor would call a non-distorted state---no active psychosis, no acute crisis, no transient situational response. I have a formed, persistent, reflectively endorsed position: competent adults with genuine decisional capacity should be permitted to choose death. I hold that position on a Tuesday when everything's fine, when the farm smells good and the conversations with my sister are gold and the work is meaningful. The argument does not require me to be suffering. It requires me to be honest.

That's the only credential I've ever claimed.

% --- Section II ---
\section{What Actually Happened at Eighteen}

\subsection{The Inventory}

250 amitriptyline. Two boxes of Benadryl. Some Percocet. Most of a bottle of Stoli. Whatever else was rattling around in the medicine cabinet that looked useful. I dumped it into a mixing bowl and ate it like fucking Cap'n Crunch, washed it down with the vodka and milk. I love milk.  No tears. No ambivalence I can locate in retrospect. Just the flat, quiet sense of arriving at the end of something that had been going on too long.

There was no drama in it. That matters. The popular narrative of suicide requires spectacle---performance of suffering, final communications, a crisis arc with recognizable beats. Mine had none. It was administrative. A task I had decided to complete, and then completed, with the same internal register as doing the dishes.

I understood the consequences: death, brain damage, vegetative continuation. I understood the alternatives: more therapy, different medication, time, whatever. I had considered and declined them. Not in a fog. Not in a panic. In the specific quiet that arrives when a person has genuinely made up their mind.

I chose Door \#1.

\subsection{The Death, Briefly}

They coded me. I was, by all clinical definitions, dead. What was present in that interval is difficult to report without sounding like a Netflix pitch. It wasn't religious. There was no tunnel, no light, no parade of deceased relatives. There was an overwhelming quality of \textit{home}---the specific sensation of having returned somewhere, of a long and exhausting trip finally, mercifully over. Peace is too small a word. Peace implies the previous absence of disturbance. This was more total than that. This was the absence of the question.

Then the universe ran in reverse. The specific sickening knowledge---before my eyes opened, before I could name it---that I was being brought back. I hated it immediately and completely. Not the pain. There wasn't much pain. The violation. The unauthorized reversal of a decision I had made.

\subsection{Three Days and a Promise}

Psych ward. Three days. Institutional kindness performed in an institutional key---the specific variety of care that is mostly about ensuring the paperwork survives a lawsuit. I don't begrudge the individual clinicians; they were working within a system built not around the patient's interior life but around the facility's liability exposure. That's the system. That's what we built.

They convened the people who loved me and made me apologize to each of them, watch them cry, absorb the damage my choice had done to people I cared about. I said the script: life is hard, I was sad, I'm sorry, you're enough.

The last part was even true. They \textit{were} enough---are enough, still. That's the piece the response never makes room for: the choice wasn't a verdict on them. It wasn't a referendum on whether love is real or whether the people around me had failed. I knew they hadn't. I was simply done. Those can coexist. The system cannot tolerate that coexistence, so it pretends the apology corrects the logic. It doesn't. It just silences it temporarily.

Then they discharged me with a fresh prescription for the exact same amitriptyline that almost did the job.

Read that again. Revive me. Extract a promise. Hand me the reload.

This is not a system oriented toward outcomes. It is a system oriented toward clean paperwork and dischargeable patients. The promise I made was extracted under the specific duress of watching people I love cry---relational coercion, perfectly effective, ethically indistinguishable from leverage. I honored it. Still am. But I never agreed it was right. I agreed because I love those people, and the state had engineered a situation in which honoring love and accepting the state's claim on my body were presented as the same gesture. They are not.

% --- Section III ---
\section{The Interminable Interim}

\subsection{Living on Borrowed Consent}

For years after, I was here on someone else's terms. Not because I had converted to the view that my presence was obligatory, not because the psych ward's arguments had landed---they hadn't---but because I had made a promise to specific people I loved and I was not going to break it. This is a distinction the protectionist model cannot accommodate: I was alive not because I had been corrected, but because I had been obligated. These are not the same thing, and conflating them is precisely how the system flatters itself into thinking it works.

Living on borrowed consent is its own particular weather. You inhabit the days. You do the things. You find genuine pockets of interest---the farm, the work, the conversations---while carrying in the background the knowledge that you're here on a pinky-swear and not because you've resolved the underlying question. It's not unlivable. It's not even, most of the time, actively miserable. It's more like being a long-term guest in a house you didn't choose, making the best of the amenities.

\subsection{The Books}

I wrote two books during this period, both of which were attempts to get underneath the problem---to find the mechanism and either dismantle it or understand it well enough to carry it without screaming.

\textit{The Serpent's Sentence} was about language. Specifically: how the linguistic construction of a narrating self creates an exile from immediate experience. The ``I'' that speaks is always slightly after the moment---a grammatical artifact that describes experience rather than inhabiting it---and the chronic sensation of being outside one's own life, watching it through glass, may not be a symptom of pathology. It may be a structural consequence of the way language builds consciousness. The question the book was trying to answer: if the narrator is always an artifact, can you get underneath it? Can you get back to the thing the narrator is narrating?

\textit{The Mathematics of Connection} came from somewhere more specific. Sam went to prison. The question it forced---more urgently, more nakedly than anything I'd theorized about---was whether genuine connection persists when all normal channels fail. Whether the link between two people is a function of the medium, or something prior to medium. The grief of that absence pushed the question out of philosophy and into something more like physics: either the connection is real and substrate-independent, or it isn't and everything is just medium.

I was reverse-engineering done. Either the mechanism would reveal itself as dismantlable---grief transforming into something livable, the hum gradually fading---or it would become understood well enough to coexist with.

Eight months ago, approximately, the hum stopped. I did not notice in real time. I noticed retroactively, when I went looking for it and found crickets.

\subsection{What Changed and What Didn't}

I like things now. The farm and my life there is meaningful. Conversations with my sister have a quality they didn't have when I was carrying the hum---more present, less performed. The work matters to me intrinsically without me having to argue myself into believing it. I am choosing to be here, genuinely, daily.

That's new. That's real.

But the realization changed exactly nothing about my politics.

Choice implies the option to un-choose. That's not a threat or a plan. It's the logical structure of voluntary action. If I'm here by choice, then not-being-here was also a choice that was available, and the availability of that choice is precisely what makes the choosing meaningful. A life lived under conscription is not the same as a life freely inhabited, even if the behaviors look identical from the outside.

Capability is not conscription. My existence is not a debt I owe the universe, my family, my community, or whatever story they're telling about who I am in their lives. The fact that I pay into those stories willingly does not mean the payment was mandatory. Mandatory contributions have no moral weight. Voluntary ones do. I want the weight.

% --- Section IV ---
\section{The Real Line: Capacity, Not Suffering}

\subsection{Why Suffering Is the Wrong Criterion}

The dominant MAiD frameworks predicate eligibility on suffering: unbearable, irremediable, demonstrable suffering. The intuition behind this is legible. We want to limit exit rights to people who are genuinely in extremis---people for whom continued existence is clearly a burden rather than a resource. The suffering criterion is an attempt to operationalize that intuition.

It fails in at least three directions.

First: suffering is not the same as impaired capacity. A person in profound pain may retain perfect decisional capacity---clear understanding of consequences, stable preferences, full comprehension of alternatives---while a person experiencing minimal subjective suffering may have thoroughly compromised capacity due to acute psychosis. The two dimensions are independent. Using one to gate the other is a category error.

Second: the suffering criterion encodes a paternalistic judgment about which reasons for exit are acceptable. It says: you must be broken enough, in a legible enough way, before we will countenance your decision. A person who has lived eighty years, completed everything they came to complete, and simply does not wish to continue---without unbearable pain, without terminal illness---has an equally valid claim. Their reasons may be harder for us to categorize. That is our problem, not theirs.

Third: ``unbearable'' is not a clinical measurement. It's a self-report. We accept the self-report in terminal cases because the physical substrate corroborates it. For psychiatric or existential cases, we treat the self-report with suspicion because the substrate is interior and invisible. This asymmetry is not principled. It is a bias toward the legible.

Suffering is a proxy for something we actually care about: genuine autonomous preference. Capacity is the direct measure. Prefer the direct measure.

\subsection{The Two-Stage Gate}

This framework proposes two stages of assessment: not gatekeeping in the pejorative sense, but genuine epistemic work---distinguishing authentic, stable, capacitated preference from transient crisis, distorted judgment, or external coercion.

\textbf{Stage 1: Decisional capacity itself.} The individual must demonstrate understanding of the permanence and irreversibility of death; documented awareness of available alternatives and their realistic outcomes; absence of acute distorting states; and persistence of the preference over time.

Awareness of alternatives is the threshold, not utilization. Informed refusal is the standard medicine applies everywhere else. You may decline chemotherapy without having sat through a single round, provided you understand what it offers and what declining forfeits. This framework applies the same standard to psychiatric and existential cases. The assessment question is not ``have you tried an SSRI'' but ``tell me what you understand about SSRIs and what led you to decline them.'' Comprehension is the gate. Compliance is not.

\textbf{Stage 2: Metacognitive capacity.} Can the individual understand \textit{why} the assessment process exists? Can they engage with safeguards as rational protections rather than obstacles to rage at? Can they tolerate a waiting period---not because waiting changes anything, but because the willingness to wait is itself evidence of the resolve being assessed?

The relationship between the stages is diagnostic. Failure of Stage 2 while claiming Stage 1 is a signal: genuine capacity includes understanding why we don't hand barbiturates to someone mid-meltdown. A person who grasps their own situation well enough to claim autonomy should be able to grasp why the process exists. If they can't, the Stage 1 claim deserves scrutiny.

The system responded by violating me, extracting a promise under duress, and handing me the ammunition for Round 2. That is not a coherent policy response to a capacitated autonomous choice. It is a liability ritual. We can do better.

% --- Section V ---
\section{``But External Pressures!'' Cry the Paternalists}

\subsection{The Objection, Stated Fairly}

The most substantive objection to capacity-based exit rights is the external pressures concern. The argument: if we permit MAiD for anyone with intact capacity, we will effectively permit it for people who are choosing death because of poverty, isolation, inadequate housing, or the internalized sense that they are a burden. Society fails people in these situations. Legalizing exit for capacitated individuals means legalizing exit for people whose capacities are intact but whose options have been structurally eliminated. We will call it autonomy while executing people for being poor.

This is not a stupid argument. It deserves a straight answer.

\subsection{The Straight Answer}

The argument proves too much. Any time we respect a choice made under constrained circumstances, we face the same problem: is this the choice they would make with better options? The person who takes a dangerous job because it's the only available income---is their consent to occupational hazard genuine? If the answer is ``no choice made under unjust constraints is truly free, so we must override all such choices,'' then we have endorsed a paternalism so comprehensive that autonomy becomes a fiction reserved for the materially comfortable.

This is not a theory of freedom. It is a theory of guardianship, applied selectively to people whose circumstances make their choices politically inconvenient.

The correct response to a capacitated person choosing death because poverty makes their life intolerable is not to deny them the choice. It is to stop making people poor. The external pressure objection is most powerful not as a reason to restrict exit rights but as a reason to deliver what we've been promising: housing, income, care, the basic material infrastructure of a livable life.

Dismissing a choice because the reasons are ``external'' rather than ``internal'' is a distinction without a principled foundation. If capacity holds, external pressures don't veto self-ownership. Fuck right off with ranking whose suffering qualifies.

\subsection{The Forcing Function}

Here is the political argument that paternalists are trying not to see: a genuine right to exit creates an obligation to make staying worth it.

Canada's numbers already show this mechanism in early form. Track 2 recipients cite housing instability, inadequate disability supports, and social isolation as primary drivers at a rate that is not incidental or anecdotal. If the state must record, on official death documentation, that a citizen chose MAiD because they had no housing and no supports---that is a different kind of public record than a death certificate reading ``complications of chronic illness.'' The specificity is politically actionable. It names the failure. It creates a constituency with standing.

Expand eligibility to the capacitated non-suffering population, and that mechanism becomes more powerful. Every documented case citing poverty or abandonment becomes a line item in the state's record of its own failures. You cannot warehouse people in undignified conditions and deny them the option to leave without eventually facing the arithmetic: the door must open, or the room must improve. This framework insists on both. Dignity in staying, or dignity in leaving. You don't get to deliver neither and call it compassion.

% --- Section VI ---
\section{How to Do This Without Being Monsters}

\subsection{Process as Filter}

The procedural concern---that capacity-based eligibility will funnel impulsive or crisis-driven individuals toward irreversible decisions---is legitimate and addressable. The process itself is the primary filter.

Anyone willing to engage with a multi-stage, multi-month, multiply-witnessed clinical assessment process is already expressing a quality of resolve that distinguishes them from impulsive crisis. People in acute crisis do not tolerate bureaucracy with equanimity. The affective urgency of crisis is incompatible with the slow grind of documented assessment. This isn't a cynical observation---it's a structural feature of the distinction we're trying to draw. Use it.

Impulsive cases grab whatever's available. Resolved cases will wait, document, engage, revoke if they change their minds, and re-engage. The truly capacitated don't flinch at calendars. They might find the waiting tedious. They will tolerate it. The process itself is evidence of resolution.

\subsection{What Must Happen}

Every realistic alternative must be presented---not as a prerequisite to clear before eligibility is granted, but as information to which the individual is entitled. Housing, income support, psychiatric care, palliative options, disability resources, peer support: the full buffet, discussed with sufficient depth for genuine comprehension. Referrals where the individual is open to them. No forced trials. If the person has understood the alternatives and declined them, that declination must be respected. The alternative is a system that treats options as mandatory prerequisites---which is precisely the double standard that has made psychiatric MAiD eligibility politically impossible for decades. No forced ``try this for six months'' horseshit.

Spiritual and existential consultation should be offered once, clearly and without pressure. If the person wants it, provide it. If they decline, move on. Repeated offers shade into proselytizing; proselytizing disguised as clinical care is its own coercion.

The process must be internally consistent from first assessment to final provision: multiple independent evaluations, documented at each stage; alternatives offered and responses recorded; capacity re-assessed at intervals as due diligence, not suspicion. Revocation must be available at any point---including the moment before provision. No ``deathday locked in'' procedure. Active, ongoing consent, verifiable until the end.

Waiting periods: the honest answer is that I don't know the magic number, and neither does anyone who claims certainty. Belgium runs approximately one month for non-terminal cases. Canada's Track 2 floor is ninety days. The Netherlands applies minimal waiting requirements. All of these seem defensible for people whose resolve is genuine; none will deter the truly decided. The waiting is not punitive. It is diagnostic. If resolve is iron, time just becomes preparation, anticipation, peace. The calendar is irrelevant to the resolved.

And yeah---I've been dead. It's fine. You're gone. It doesn't hurt. No notes from the other side complaining.

% --- Section VII ---
\section{Why My Argument Hits Different Now}

\subsection{The Credibility of the Non-Suicidal Advocate}

I am not suicidal. I'm stable, choosing, contributing, functioning in every dimension the system uses to measure whether a person is okay. I find satisfaction in the farm, in the work, in the conversations that matter. I am choosing this particular shitshow daily with something resembling enthusiasm.

And I still say: my consciousness is mine. Not society's asset. Not my family's property. Not the state's portfolio.

Arguing this from stability is not paradoxical. It is the only position from which the argument is fully credible. The person making this case from within crisis can be dismissed---not because they're wrong, but because the dismissal is always available: \textit{they're not thinking clearly, they're in pain, they'd feel differently later}. I can't be dismissed that way. I am thinking clearly. I am not in pain. And I've spent twenty years finding out what I'd feel like later.

This is what later looks like. The argument survived.

\subsection{The Mistakes Are Mine to Make}

The objection I find most emotionally resonant, even when I think it's philosophically weak, is the regret argument: what if I had been permitted to exit at eighteen, and I would have been happy at thirty-eight? Doesn't that possibility justify forced continuation?

My answer: yes. And?

I might have made a mistake. The decision, made at eighteen with genuine capacity, might have foreclosed a future that was worth having. This is also true of every irreversible decision any person makes under any framework of autonomy. The person who declines a medical treatment and dies from a treatable condition might have been happy if they'd complied. We do not, as a general principle, use the possibility of regret-that-cannot-be-experienced as grounds for overriding present autonomous choice.

The regret argument is a special pleading reserved for exit decisions because we find them morally and aesthetically disturbing. The disturbance is real. It is not a principle. My mistake to make---if it would have been a mistake---is mine. That's what ownership means.

I stayed for years because people I love asked me to. I'm glad I did. That gladness does not retroactively make the coercion correct. It just means the outcome was good. Good outcomes don't justify any method of achieving them. We know this. We believe it about everything else.

% --- Section VIII ---
\section{Final Landing}

I don't want to die. That is genuinely, weirdly, wonderfully new. I am here by choice. I am choosing to keep being here. Every day that I do is a voluntary act, and that voluntariness is the only thing that gives it any moral weight whatsoever.

The world does not own my exit because it is uncomfortable watching me leave. Some people get done. If they are lucid---if they have demonstrated stable, informed, persistent preference, if they understand the stakes and have weighed the alternatives and arrived at their answer with clear eyes---let them go.

Not because death is better than life. Not because pain automatically qualifies anyone. Not because suffering is sufficient. Because capacity equals sovereignty. Because the only condition under which presence means anything is the condition under which it is genuinely chosen. And the only condition under which it can be genuinely chosen is the condition under which not-choosing it was also available.

Bar the exit and you don't save a life. You take ownership of one.

Unlock it, and everyone who stays does so as an act of will. That's the only version of a life that has moral weight. That's the only version worth anything.

Replace sanctity-of-life dogma with the sovereignty of the living. Make the room worth staying in. And if you won't---if the room is genuinely uninhabitable and you won't fix it---then at minimum have the decency to unlock the door.

My argument does not need me to be miserable to be true. It just needs to be honest.

It is.

\bigskip
\begin{flushleft}
\textbf{References}\\[0.5em]
\small
Health Canada. (2025). \textit{Sixth Annual Report on Medical Assistance in Dying in Canada} (2024 data). Ottawa: Health Canada.\\[0.3em]
Health Canada. (2023--2024). \textit{Track 2 case documentation and socioeconomic analyses}. Ottawa: Health Canada.\\[0.3em]
Ontario Office of the Chief Coroner. (2022--2025). \textit{Reviews of MAiD cases involving socioeconomic drivers}. Toronto.\\[0.3em]
United Nations Special Rapporteur on the Rights of Persons with Disabilities. (2019). \textit{Report on autonomy, legal capacity, and MAiD}. Geneva: UN Human Rights Council.\\[0.3em]
Bogner, J. (2024). \textit{The Serpent's Sentence}.\\[0.3em]
Bogner, J. (2025). \textit{The Mathematics of Connection}.
\end{flushleft}

\end{document}
